% ============================================================
%  conclusao.tex  –  Conclusão: GMM-Bayes vs. classificadores
%  clássicos nos cinco datasets do Trabalho 5.
%  Inclua no documento principal com 
\chapter{Resultados}

 
Este trabalho avaliou o classificador Bayesiano com Misturas de Gaussianas (GMM-Bayes,
$J=3$ componentes por classe) em cinco conjuntos de dados, comparando-o com LDA, QDA,
KNN ($k=5$) e DMC ao longo de 20 realizações independentes com divisão estratificada
80\%/20\%. A Tabela~\ref{tab:resumo_geral} consolida os resultados obtidos.
 
\begin{table}[htbp]
  \centering
  \caption{Resumo geral — acurácia média (\%) dos classificadores por dataset.}
  \label{tab:resumo_geral}
  \begin{tabular}{lccccc}
    \hline
    \textbf{Dataset}      & \textbf{GMM-Bayes} & \textbf{LDA}  & \textbf{QDA}  & \textbf{KNN}  & \textbf{DMC}  \\
    \hline
    Artificial I          & \textbf{100,00}    & 97,69         & 98,46         & \textbf{100,00} & 88,85       \\
    Breast Cancer         & 94,12              & 96,23         & 95,09         & \textbf{96,49}  & 93,46       \\
    Dermatology           & 89,44              & \textbf{96,94} & 30,56        & 96,32         & 96,25       \\
    Íris                  & 95,83              & \textbf{98,33} & 97,83        & 95,83         & 86,83       \\
    Coluna Vertebral      & 80,00              & 77,98         & \textbf{83,23} & 77,34        & 73,47       \\
    \hline
  \end{tabular}
\end{table}

\section{Análise por dataset}

\subsection{Artificial I.}
O GMM-Bayes obteve acurácia média perfeita (100,00\%), empatando com o KNN e superando
QDA (98,46\%) e LDA (97,69\%). 
\begin{table}[htbp]
  \centering
  \caption{Resultados detalhados --- Artificial I.}
  \label{tab:det_artificial}
  \begin{tabular}{lcccc}
    \hline
    \textbf{Classificador} & \textbf{Média (\%)} & \textbf{DP (\%)} & \textbf{Mín (\%)} & \textbf{Máx (\%)} \\
    \hline
    GMM-Bayes & 100,00 & 0,00 & 100,00 & 100,00 \\
    LDA       & 97,69  & 2,55 & 92,31  & 100,00 \\
    QDA       & 98,46  & 1,88 & 96,15  & 100,00 \\
    KNN       & 100,00 & 0,00 & 100,00 & 100,00 \\
    DMC       & 88,85  & 5,82 & 73,08  & 100,00 \\
    \hline
  \end{tabular}
\end{table}


\begin{figure}[h]
    \centering
    \includegraphics[width=0.5\textwidth]{documentos/imagens/artificial/superficies_decisao.png}
   \caption{Superfície de decisão do discriminante quadratico — Dataset Artificial I}
    \label{fig:iris_gaussians}
\end{figure}

\begin{figure}[h]
    \centering
    \includegraphics[width=0.5\textwidth]{documentos/imagens/artificial/matriz_confusao.png}
     \caption{Matrizes de confusão dos três classificadores — Dataset Artificial I}
    \label{fig:iris_gaussians}
\end{figure}

Esse resultado é esperado: o dataset foi gerado sinteticamente
com três clusters bem separados por classe, estrutura multimodal que os modelos de mistura
capturam naturalmente. A acurácia perfeita e desvio padrão nulo indicam que os três
componentes do GMM representam com exatidão a geometria dos dados, sem necessidade de
fronteiras mais complexas como as do KNN.

\subsection{Breast Cancer (WDBC).}
No \textit{Wisconsin Diagnostic Breast Cancer}, o GMM-Bayes obteve acurácia média de
\begin{table}[htbp]
  \centering
  \caption{Resultados detalhados --- Breast Cancer (WDBC).}
  \label{tab:det_breast}
  \begin{tabular}{lcccc}
    \hline
    \textbf{Classificador} & \textbf{Média (\%)} & \textbf{DP (\%)} & \textbf{Mín (\%)} & \textbf{Máx (\%)} \\
    \hline
    GMM-Bayes & 94,12 & 2,02 & 90,35 & 97,37 \\
    LDA       & 96,23 & 1,57 & 92,11 & 98,25 \\
    QDA       & 95,09 & 1,58 & 92,98 & 99,12 \\
    KNN       & 96,49 & 1,59 & 92,98 & 99,12 \\
    DMC       & 93,46 & 1,56 & 90,35 & 96,49 \\
    \hline
  \end{tabular}
\end{table}
\begin{figure}[h]
    \centering
    \includegraphics[width=0.5\textwidth]{documentos/imagens/cancer/matriz_confusao.png}
     \caption{Matrizes de confusão dos três classificadores — Dataset Cancer}
    \label{fig:iris_gaussians}
\end{figure}
A maior variabilidade do GMM-Bayes  com uma faixa de quase $7$ pontos percentuais entre a pior e a melhor
realização reflete a dificuldade de estimar $J=3$ componentes em $d=30$ dimensões;
na melhor realização (13ª), o modelo chegou a $97{,}37\%$, mas em outras rodadas caiu
para $\approx 90\%$. Em alta dimensionalidade com apenas duas classes, modelos mais
parcimoniosos como o LDA e o KNN mostraram-se mais estáveis.

\subsection{Dermatology.}
O GMM-Bayes atingiu $89{,}44\%$ de acurácia média de 
\begin{table}[htbp]
  \centering
  \caption{Resultados detalhados --- Dermatology.}
  \label{tab:det_dermatology}
  \begin{tabular}{lcccc}
    \hline
    \textbf{Classificador} & \textbf{Média (\%)} & \textbf{DP (\%)} & \textbf{Mín (\%)} & \textbf{Máx (\%)} \\
    \hline
    GMM-Bayes & 89,44 & 5,92 & 79,17 & 97,22 \\
    LDA       & 96,94 & 1,68 & 93,06 & 98,61 \\
    QDA       & 30,56 & 0,00 & 30,56 & 30,56 \\
    KNN       & 96,32 & 1,66 & 91,67 & 98,61 \\
    DMC       & 96,25 & 2,16 & 91,67 & 100,00 \\
    \hline
  \end{tabular}
\end{table}

\begin{figure}[h]
    \centering
    \includegraphics[width=0.5\textwidth]{documentos/imagens/dermatologia/matriz_confusao.png}
     \caption{Matrizes de confusão dos três classificadores — Dataset Dermatologia}
    \label{fig:iris_gaussians}
\end{figure}
equivalente à predição
da classe majoritária, evidenciando severo overfitting em $d=34$ atributos. O
GMM-Bayes, apesar de mais expressivo, apresentou alta variabilidade (amplitude de
$\approx 18$ pontos percentuais), alcançando $97{,}22\%$ na melhor realização (7ª),
mas performando abaixo dos modelos lineares e de distância na média geral.

\subsection{Íris.}
O GMM-Bayes obteve\begin{table}[htbp]
  \centering
  \caption{Resultados detalhados --- Íris.}
  \label{tab:det_iris}
  \begin{tabular}{lcccc}
    \hline
    \textbf{Classificador} & \textbf{Média (\%)} & \textbf{DP (\%)} & \textbf{Mín (\%)} & \textbf{Máx (\%)} \\
    \hline
    GMM-Bayes & 95,83 & 3,78 & 86,67 & 100,00 \\
    LDA       & 98,33 & 1,97 & 93,33 & 100,00 \\
    QDA       & 97,83 & 2,42 & 93,33 & 100,00 \\
    KNN       & 95,83 & 2,07 & 93,33 & 100,00 \\
    DMC       & 86,83 & 4,88 & 76,67 & 96,67 \\
    \hline
  \end{tabular}
\end{table}

\begin{figure}[h]
    \centering
    \includegraphics[width=0.5\textwidth]{documentos/imagens/iris/superficies_decisao.png}
   \caption{Superfície de decisão do discriminante quadratico — Dataset iris}
    \label{fig:iris_gaussians}
\end{figure}
\begin{figure}[h]
    \centering
    \includegraphics[width=0.5\textwidth]{documentos/imagens/iris/matriz_confusao.png}
     \caption{Matrizes de confusão dos três classificadores — Dataset iris}
    \label{fig:iris_gaussians}
\end{figure}
Na 3ª realização, o GMM-Bayes chegou à
acurácia máxima de $100\%$, mas sua média foi penalizada por realizações com
$\approx 86{,}67\%$, sugerindo que três componentes por classe introduziram
complexidade desnecessária para um dataset aproximadamente unimodal e linearmente
separável.

\subsection{Coluna Vertebral.}
No dataset \textit{Vertebral Column}, o GMM-Bayes alcançou $80{,}00\%$ \begin{table}[htbp]
  \centering
  \caption{Resultados detalhados --- Coluna Vertebral.}
  \label{tab:det_vertebral}
  \begin{tabular}{lcccc}
    \hline
    \textbf{Classificador} & \textbf{Média (\%)} & \textbf{DP (\%)} & \textbf{Mín (\%)} & \textbf{Máx (\%)} \\
    \hline
    GMM-Bayes & 80,00 & 3,44 & 74,19 & 88,71 \\
    LDA       & 77,98 & 4,48 & 66,13 & 87,10 \\
    QDA       & 83,23 & 3,93 & 74,19 & 88,71 \\
    KNN       & 77,34 & 3,36 & 69,35 & 83,87 \\
    DMC       & 73,47 & 5,21 & 62,90 & 82,26 \\
    \hline
  \end{tabular}
\end{table}
\begin{figure}[h]
    \centering
    \includegraphics[width=0.5\textwidth]{documentos/imagens/vertebral/superficies_decisao.png}
   \caption{Superfície de decisão do discriminante quadratico — Dataset Vertebral}
    \label{fig:iris_gaussians}
\end{figure}
\begin{figure}[h]
    \centering
    \includegraphics[width=0.5\textwidth]{documentos/imagens/vertebral/matriz_confusao.png}
     \caption{Matrizes de confusão dos três classificadores — Dataset Vertebral}
    \label{fig:iris_gaussians}
\end{figure}

Com $d=6$ atributos
e três classes, o número de parâmetros do GMM é mais controlável do que nos
datasets de alta dimensionalidade, permitindo que o modelo capture não-linearidades
moderadas. Ainda assim, o QDA  que estima uma única covariância completa por
classe  apresentou melhor média, indicando que a multimodalidade implícita em
$J=3$ não se traduziu em vantagem consistente neste conjunto.
 
\section*{Conclusões gerais}
 
Com base nos experimentos realizados e na atualização da implementação para EM explícito,
as seguintes conclusões podem ser destacadas:
 
\begin{enumerate}
 
  \item \textbf{O GMM-Bayes confirmou bom desempenho em cenários multimodais.}
  No Artificial~I, o método manteve acurácia média de 100,00\%, empatando com KNN e
  superando os demais classificadores lineares/quadráticos na média.
 
  \item \textbf{Nos dados reais, o desempenho foi competitivo, porém não dominante.}
  As médias do GMM-Bayes foram 94,12\% (Breast Cancer), 89,44\% (Dermatology),
  95,83\% (Íris) e 80,00\% (Coluna Vertebral). Em Breast Cancer e Íris, o método ficou
  próximo dos melhores; em Vertebral, ficou acima de LDA/KNN/DMC, mas abaixo de QDA.
 
  \item \textbf{A nova implementação passou a usar EM explícito no código.}
  Diferentemente da versão anterior baseada em chamada implícita de biblioteca, o código
  atual implementa diretamente o Passo E (responsabilidades) e o Passo M (atualização de
  pesos, médias e covariâncias), com inicialização por K-Means e regularização numérica.
 
  \item \textbf{Validação da versão explícita no Iris foi bem-sucedida.}
  No teste de caminho crítico executado em \texttt{gmm\_bayes.py}, as variantes obtiveram:
  Full = 95,56\%, Diagonal = 93,33\%, Spherical = 95,56\% e Tied = 97,78\%, com melhor
  resultado para a configuração \emph{tied}.
 
  \item \textbf{A escolha de modelagem continua crítica para generalização.}
  Mesmo com EM explícito, os resultados reforçam que desempenho depende de dimensionalidade,
  separabilidade e número de componentes por classe, mantendo válida a recomendação de
  seleção criteriosa de hiperparâmetros (por exemplo, via validação cruzada/BIC/AIC).
 
\end{enumerate}
 
Em síntese, o GMM-Bayes com EM explícito mostrou-se uma alternativa tecnicamente consistente
e interpretável para classificação supervisionada probabilística. O método foi robusto nos
testes executados e permanece especialmente atrativo quando há evidência de multimodalidade
intra-classe, embora modelos mais parcimoniosos possam generalizar melhor em alguns cenários
de alta dimensionalidade.

% ============================================================

\section{Conclusão}

Este trabalho avaliou o classificador Bayesiano com Misturas de Gaussianas (GMM-Bayes,
$J=3$ componentes por classe) em cinco conjuntos de dados, comparando-o com LDA, QDA,
KNN ($k=5$) e DMC ao longo de 20 realizações independentes com divisão estratificada
80\%/20\%. A Tabela~\ref{tab:resumo_geral} consolida os resultados obtidos.

\begin{table}[htbp]
  \centering
  \caption{Resumo geral — acurácia média (\%) dos classificadores por dataset.}
  \label{tab:resumo_geral}
  \begin{tabular}{lccccc}
    \hline
    \textbf{Dataset}      & \textbf{GMM-Bayes} & \textbf{LDA}  & \textbf{QDA}  & \textbf{KNN}  & \textbf{DMC}  \\
    \hline
    Artificial I          & \textbf{100,00}    & 97,69         & 98,46         & \textbf{100,00} & 88,85       \\
    Breast Cancer         & 94,12              & 96,23         & 95,09         & \textbf{96,49}  & 93,46       \\
    Dermatology           & 89,44              & \textbf{96,94} & 30,56        & 96,32         & 96,25       \\
    Íris                  & 95,83              & \textbf{98,33} & 97,83        & 95,83         & 86,83       \\
Coluna Vertebral      & 80,00              & 77,98         & 83,23         & 77,34        & 73,47        \\
    \hline
  \end{tabular}
\end{table}

\subsection*{Análise por dataset}

\paragraph{Artificial I.}
O GMM-Bayes obteve acurácia média perfeita (100,00\%), empatando com o KNN e superando
QDA (98,46\%) e LDA (97,69\%). Esse resultado é esperado: o dataset foi gerado sinteticamente
com três clusters bem separados por classe, estrutura multimodal que os modelos de mistura
capturam naturalmente. A acurácia perfeita e desvio padrão nulo indicam que os três
componentes do GMM representam com exatidão a geometria dos dados, sem necessidade de
fronteiras mais complexas como as do KNN.

\paragraph{Breast Cancer (WDBC).}
Com 30 atributos e apenas 2 classes, o GMM-Bayes (94,12\%) ficou abaixo do LDA (96,23\%)
e do KNN (96,49\%). Em alta dimensionalidade com poucas classes, estimar $J \cdot d(d+1)/2$
parâmetros de covariância por componente torna-se instável, favorecendo modelos mais
parcimoniosos como o LDA. O QDA com regularização também superou o GMM-Bayes (95,09\%),
reforçando que a hipótese de gaussiana única por classe, com pooling ou regularização, é mais
adequada para este dataset.

\paragraph{Dermatology.}
O GMM-Bayes atingiu 89,44\%, um resultado moderado frente ao LDA (96,94\%), KNN (96,32\%)
e DMC (96,25\%). O resultado mais crítico neste dataset foi o colapso do QDA padrão (30,56\%),
que corresponde à acurácia de um classificador que prediz apenas a classe majoritária,
evidenciando severo overfitting com 34 atributos e poucas amostras por classe. O GMM-Bayes,
embora mais expressivo que o LDA, também sofreu com a alta dimensionalidade, não sendo capaz
de superar os modelos que impõem restrições mais fortes à covariância. Neste cenário, modelos
lineares ou de distância mostraram-se mais robustos.

\paragraph{Íris.}
O GMM-Bayes (95,83\%) ficou abaixo do LDA (98,33\%) e do QDA (97,83\%), empatando com
o KNN. O dataset Íris é amplamente reconhecido como aproximadamente linearmente separável
após a remoção da classe \textit{setosa}, o que favorece discriminantes lineares. A adoção de
três componentes por classe, possivelmente com sobreposição dentro de cada classe unimodal,
introduziu uma complexidade desnecessária que penalizou levemente o GMM-Bayes em relação
ao LDA.

\paragraph{Coluna Vertebral.}
Os resultados do dataset Coluna Vertebral ainda serão incorporados posteriormente.

\subsection*{Conclusões gerais}

Com base nos experimentos realizados, as seguintes conclusões podem ser extraídas:

\begin{enumerate}

  \item \textbf{GMM-Bayes é mais eficaz quando os dados são genuinamente multimodais.}
  No Artificial~I, onde cada classe foi gerada com múltiplos focos gaussianos, o modelo
  atingiu desempenho perfeito. A capacidade de representar distribuições multimodais
  é a principal vantagem dos GMMs sobre os classificadores de componente único
  \cite{bishop2006,theodoridis2009}.

  \item \textbf{Alta dimensionalidade penaliza o GMM-Bayes.}
  Nos datasets Breast Cancer (30 atributos) e Dermatology (34 atributos), o GMM-Bayes
  ficou abaixo dos modelos mais simples. A estimação de $J \cdot d(d+1)/2$ parâmetros de
  covariância por componente exige um número de amostras por classe muito maior do que o
  disponível, tornando o modelo suscetível a overfitting \cite{duda2001}.

  \item \textbf{LDA mantém robustez em cenários de alta dimensão e separação linear.}
  Para Íris, Breast Cancer e Dermatology, o LDA apresentou os melhores ou segundo melhores
  resultados entre os classificadores Bayesianos. Ao impor uma covariância pooled comum,
  reduz drasticamente o número de parâmetros e generaliza melhor com amostras limitadas
  \cite{duda2001,bishop2006}.

  \item \textbf{O colapso do QDA no Dermatology é um alerta sobre overfitting.}
  A acurácia de 30,56\% do QDA padrão — equivalente à predição da classe majoritária —
  demonstra que, sem regularização, estimar uma matriz de covariância completa por classe
  em dimensão $d=34$ com poucas amostras é numericamente instável. Essa instabilidade
  justifica o uso de regularização de Tikhonov ou interpolação com a covariância pooled
  \cite{bishop2006}.

  \item \textbf{A escolha do número de componentes $J$ é crítica.}
  Usar $J=3$ foi adequado para o Artificial~I mas potencialmente excessivo para datasets
  com classes unimodais (Íris) ou de alta dimensão (Dermatology, Breast Cancer). Uma
  seleção de $J$ por validação cruzada, como proposto na metodologia, é necessária para
  equilibrar expressividade e generalização \cite{theodoridis2009}.

  \item \textbf{O DMC apresentou o pior desempenho consistente.}
  Em todos os datasets, o DMC ficou na última posição ou próximo dela, confirmando sua
  limitação a classes com geometria esférica e bem separadas. A hipótese
  $\boldsymbol{\Sigma}_i = \sigma^2 \mathbf{I}$ é excessivamente restritiva para dados
  reais \cite{duda2001}.

\end{enumerate}

Em síntese, o GMM-Bayes representa um passo natural além dos classificadores Bayesianos
gaussianos clássicos, oferecendo maior flexibilidade para representar densidades multimodais.
Contudo, essa flexibilidade tem custo: demanda mais dados por classe e um procedimento
cuidadoso de seleção de hiperparâmetros. Em cenários de baixa dimensão com dados multimodais,
é competitivo ou superior a QDA e LDA; em alta dimensão, regularização e modelos mais
parcimoniosos tendem a generalizar melhor.
